\printedpage{14}
% head: A. Grothendieck : Dérivateurs, Ch. I, Édition des Universités de Montpellier II et Paris VII
% foot: Version 8 juillet 2001

and $\Sigma^\circ{}_I$ is deduced from $\Sigma_{I^\circ}$, whence
\[
(\mathcal{M}^\circ)_I(\Sigma^\circ{}_I)^{-1} \simeq (\mathcal{M}_{I^\circ}\Sigma_{I^\circ}^{-1})^\circ = \D(I^\circ)^\circ \quad .
\]

\msp{page 18}

\textbf{Philosophy.} Each time we have a functor
\[
\mathcal{A} \xrightarrow{u_0} \mathcal{A}'
\]
between two categories which can be naturally interpreted as the fundamental categories of two derivators $\D$, $\D'$, the question arises of interpreting $u_0$ as coming from a morphism of derivators
\[
\D \xrightarrow{u} \D' \quad .
\]
The latter is, in practice, always canonically determined, up to canonical isomorphism. (But I nevertheless doubt that the restriction functor to the fundamental categories
\[
\uHom(\D,\D') \longrightarrow \uHom(\mathcal{A},\mathcal{A}')
\]
is fully faithful \dots\ )

We are going to apply this philosophy to the functors $f^*$, $f_!$, $f_*$ (and possibly $f^!$, $f^?$), for $f : I \to J$, thus to the functors between $\D(I)$ and $\D(J)$. This therefore leads us to regard $\D(I)$ and $\D(J)$ as fundamental categories of suitable derivators.

\textbf{Induced derivator} by a derivator $\D$, on an index category $I \in \Dia$. It is a derivator $\D_I$ defined by
\[
\D_I(J) = \D(I \times J) \quad .
\]
The functor associated with $\D_I$, for two categories $J$, $J' \in \Dia$

\msp{page 19}
\[
\uHom(J,J')^\circ \xrightarrow{(\D_I)_{J,J'}} \uHom(\D_I(J'),\D_I(J)) = \uHom(\D(I\times J'),\D(I\times J))
\]
is the composite
\[
\uHom(J,J')^\circ \longrightarrow \uHom(I\times J, I\times J')^\circ \xrightarrow{\D_{I\times J,I\times J'}} \uHom(\D(I\times J'),\D(I\times J)) \quad ,
\]
where the first arrow is the functor
\[
f \longmapsto \id_I \times f \quad .
\]
In other words, $\D_I$ is the composite
\[
\begin{aligned}
&\Dia^\circ \longrightarrow \Dia^\circ \xrightarrow{\D} \Cat \\
&J \longmapsto I \times J \qquad .
\end{aligned}
\]
We have
\[
f^*_{\D_I} = (\id \times f)^*_{\D} \quad .
\]
If the $f_*$ (resp. the $f_!$) exist for $\D$, they exist for $\D_I$. (Likewise for every kind of multiadjoints.) Since $J \mapsto I \times J$ takes closed immersion (resp. open) to likewise, axiom \textbf{Der}$_2$ (more precisely, each of these ingredients, \textit{cf.} marginal annotations p.~7) is stable under passage from $\D$ to $\D_I$ . Likewise for \textbf{Der}$_3$. (For each new axiom we bring out, we will have to check stability.)

\printedpage{15}
% head: A. Grothendieck : Dérivateurs, Ch. I, Édition des Universités de Montpellier II et Paris VII
% foot: Version 8 juillet 2001

\msp{page 20}

Let
\[
f : I \longrightarrow I' \quad .
\]
Consider between $\D(I)$, $\D(I')$ the functors
\[
\begin{tikzcd}[column sep=huge]
\D(I) \arrow[r, shift left=3ex, "f_!"] \arrow[r, shift right=3ex, "f_*"'] & \D(I') \arrow[l, "f^*"']
\end{tikzcd}
\]
(subject to the existence of the $f_*$ resp. $f_!$). Can we interpret them in terms of morphisms of derivators
\[
\begin{tikzcd}[column sep=huge]
\D_I \arrow[r, shift left=3ex, "\overline{f_!}"] \arrow[r, shift right=3ex, "\overline{f_*}"'] & \D_{I'} \arrow[l, "\overline{f^*}"']
\end{tikzcd}
\qquad ?
\]
So we must define, for $J$ in $\Dia$, corresponding functors between $\D_I(J)$ and $\D_{I'}(J)$, \textit{i.e.} between $\D(I \times J)$ and $\D(I' \times J)$. Now we do have
\[
f \times \id_J : I \times J \longrightarrow I' \times J
\]
and we will set (subject to the existence of the last two)
\[
\left\{
\begin{aligned}
&\overline{f^*}(J) = (f \times \id_J)^*_{\D} \\[1ex]
&\overline{f_!}(J) = (f \times \id_J)^{\D}_! \\[1ex]
&\overline{f_*}(J) = (f \times \id_J)^{\D}_* \quad .
\end{aligned}
\right.
\]

\msp{page 21}

In other words, we can regard
\[
\Dia \ni I \longmapsto \D_I \quad , \qquad \D \in \DER \quad ,
\]
as functorial in $I$ in three ways (one contravariant, the other two covariant), according as we use the $f^*$, the $f_!$ (if they exist) or the $f_*$ (if they exist). We will denote by
\[
f^{\D}_! \quad , \qquad f^*_{\D} \quad , \qquad f^{\D}_*
\]
indifferently the functors between $\D(I)$ and $\D(I')$ or (more finely) between $\D_I$ and $\D_{I'}$, from which the preceding ones are deduced by passing to the corresponding fundamental categories. Likewise for $f^!$ or for $f^?$ subject to existence. For example, if we assume that the $f^!$ exist for $\D$, for every $f$ a closed immersion, then for such an $f$ we find a morphism of derivators
\[
\D_I \xleftarrow{f^!_{\D}} \D_{I'} \quad .
\]
We must see whether all these morphisms of derivators are exact, \textit{i.e.} whether for
\[
g : J \longrightarrow J'
\]
in $\Dia$, they commute with $g_!$, $g_*$.

\printedpage{16}
% head: A. Grothendieck : Dérivateurs, Ch. I, Édition des Universités de Montpellier II et Paris VII
% foot: Version 8 juillet 2001

\msp{page 22 vide}

\msp{page 23}

To make this question explicit, we are led to look at the diagram of categories
\[
\begin{tikzcd}[column sep=huge, row sep=huge]
I \times J \arrow[r, "G = \id_I \times g"] \arrow[d, "F = f \times \id_J"'] & I \times J' \arrow[d, "F' = f \times \id_{J'}"] \\
I' \times J \arrow[r, "G' = \id_{I'} \times g"] & I' \times J'
\end{tikzcd}
\qquad ,
\]
which is moreover a cartesian square in $\Cat$. The commutativity of $f^*_{\D}$ with the $g_*$ is then made explicit by the commutativity in
\[
\begin{tikzcd}[column sep=huge, row sep=huge]
\D(I \times J) \arrow[r, "G_*", ""{name=T, below, near end}] & \D(I \times J') \\
\D(I' \times J) \arrow[u, "F^*"] \arrow[r, "G'_*"] & \D(I' \times J') \arrow[u, "{F'}^*"', ""{name=R, left}]
\arrow[dashed, from=R, to=T, "\text{iso?}"]
\end{tikzcd}
\qquad .
\]
This commutativity will be ensured by later axioms on derivators. (cf. further % [sic] ultérieures -> ultérieurs, accord avec « axiomes »
below on page 59). Thus, $f^*_{\D}$ is right exact \texttt{[rather left]}, and dually, it must also be left exact \texttt{[rather right]}. So $f^*_D : \D_{I'} \to \D_I$ is \textit{exact.} Consider $f^{\D}_*$, % [sic] f^*_D -> f^*_{\mathbb{D}}, indice imprimé en D italique au lieu de \mathbb{D}
this time it is the commutativity of
\[
\begin{tikzcd}[column sep=huge, row sep=huge]
\D(I \times J) \arrow[r, "G_*"] \arrow[d, "F_*"'] & \D(I \times J') \arrow[d, "F'_*"] \\
\D(I' \times J) \arrow[r, "G'_*"] & \D(I' \times J')
\end{tikzcd}
\qquad ,
\]
it is trivial (up to canonical isomorphism). Thus, we
\msp{page 24}
find that $f^{\D}_* : \D_I \to \D_{I'}$ is \textit{left exact}, and dually $f^{\D}_!$ is \textit{right exact.} But is $f^{\D}_*$ also \textit{right} exact? This therefore means commutativity of
\[
\begin{tikzcd}[column sep=huge, row sep=huge]
\D(I \times J) \arrow[r, "G_!"] \arrow[d, "F_*"'] & \D(I \times J') \arrow[d, "F'_*", ""{name=R2, left}] \\
\D(I' \times J) \arrow[r, "G'_!", ""{name=B, above, near end}] & \D(I' \times J')
\arrow[dashed, from=B, to=R2]
\end{tikzcd}
\qquad .
\]

\printedpage{17}
% head: A. Grothendieck : Dérivateurs, Ch. I, Édition des Universités de Montpellier II et Paris VII
% foot: Version 8 juillet 2001

We do manage to define
\[
G'_!F_* \xrightarrow{?} F'_*G_!
\]
or what comes to the same thing by adjunction
\[
F_* \xrightarrow{?} {G'}^*F'_*G_!
\]
using the base change morphism
\[
{G'}^*F'_* \xrightarrow{\overset{\alpha}{\sim}} F_*G^*
\]
which is \texttt{[an]} \textit{isomorphism} (we have seen this), whence
\[
\begin{tikzcd}[column sep=large, row sep=normal]
& {G'}^*F'_*G_! \arrow[dd, "\wr"', "\alpha * G_!"] \\
F_* \arrow[ur, "?"] \arrow[dr, "F_* * \beta"'] & \\
& F_*G^*G_!
\end{tikzcd}
\]
where $\beta : \id \to G^*G_!$ is the adjunction morphism, whence the desired morphism. (We can also look for ${F'}^*(G'_!F_*) \to G_!$, and use the commutation $G_!F^* \overset{\sim}{\to} {F'}^*G'_!$, \textit{i.e.} the \textit{right} exactness of $f^*_{\D}$.)

\msp{page 25}

But is this morphism an isomorphism? I think that in the additive case it may be reasonable, but in the nonadditive case there is no reason for it to be true. We see in any case that if for every $f : I \to I'$ in $\Dia$, $f^{\D}_* : \D_I \to \D_{I'}$ is (not only left exact but also) right exact, then for every $g : J \to J'$ in $\Dia$, $g^{\D}_!$ is (not only right exact but even) left exact. % [sic] exacte -> exact, accord avec « $g^{\D}_!$ est (non seulement exact\dots) »

Let us take the derivator of presheaves of sets on variable $I$'s, its fundamental category is $\Ens$. It is a nonpointed derivator. Consider $I \to I'$ a nonsurjective closed immersion, $f_* : \D(I) \to \D(I')$. If $f^{\D}_*$ were right exact $f_*$ would have to take initial object of $\D(I)$ to initial object of $\D(I')$, unless I am mistaken (\textit{cf.} below). But this is not so, since for every $F \in \sIm\bigl(f^{\D}_* : \D(I) \to \D(I')\bigr)$ the restriction of $F$ to $I' \setminus I$ is the final sheaf $e_U = (i \mapsto e)$, and not the initial sheaf $\emptyset_U = (i \mapsto \emptyset)$.

We are going to prove this:

\textbf{Proposition.} \textit{If $\D$ is a left (right) derivator, then finite products (finite sums) exist in the $\D(I)$. Every morphism $\varphi : \D \to \D'$ of derivators which is left (right) exact commutes with the said products (sums), each $I$ giving $\varphi(I) : \D(I) \to \D'(I)$. In particular, it takes final (initial) object to likewise.}

\printedpage{18}
% head: A. Grothendieck : Dérivateurs, Ch. I, Édition des Universités de Montpellier II et Paris VII
% foot: Version 8 juillet 2001

\msp{page 26}

Saying that finite products indexed by the finite set $I$ exist in the category $\mathcal{A} = \D(e)$ (say) is equivalent to saying that the diagonal functor
\[
\mathcal{A} \xrightarrow{\Delta} \mathcal{A}^I
\]
admits a left adjoint. Now if
\[
f : I \to e
\]
is the projection of $I$ (regarded as a discrete ordered set) onto the one-point ordered set $e$, the functor $\Delta$ is identified with $f^* : \D(e) \to \D(I)$ (thanks to axiom \textbf{Der}$_2$ on the behavior of the derivator for a discrete sum of index categories). So the existence of the product amounts to that of the right adjoint $f_*$ of $f^*$, OK for left derivator. And a left exact homomorphism of left derivators, since it commutes with the $f_*$, commutes with finite products. Likewise for finite sums, and right derivators.

The preceding example uses the fact that the derivator under consideration is not pointed. But even for a pointed derivator $\D$, a morphism $f^{\D}_* : \D_I \to \D_{I'}$ is not necessarily right exact. If we take the derivator associated with an arbitrary category $\mathcal{M}$ (with $\Sigma = \emptyset$),
\[
\D_{\mathcal{M}}(I) = \uHom(I^\circ,\mathcal{M}) \overset{\text{def}}{=} \mathcal{M}(I)
\]
then the induced derivator
\[
(\D_{\mathcal{M}})_I = \{J \mapsto \D_{\mathcal{M}}(I \times J) = \uHom((I \times J)^\circ,\mathcal{M}) \simeq \uHom(J^\circ,\uHom(I^\circ,\mathcal{M}))\}
\]

\msp{page 27}

is identified with
\[
\D_{\uHom(I^\circ,\mathcal{M})} = \D_{\mathcal{M}(I)}
\]
\textit{i.e.}
\[
(\D_{\mathcal{M}})_I = \D_{\mathcal{M}(I)}
\]
N.B. Here $\mathcal{M}$ and $\D_{\mathcal{M}}$ determine each other and the functor
\[
\mathcal{M} \longmapsto \D_{\mathcal{M}}
\]
\[
\Cat \longrightarrow \DER
\]
is 2-faithful, \textit{i.e.} if $\mathcal{M}$, $\mathcal{M}'$ are two categories
\[
\uHom(\mathcal{M},\mathcal{M}') \xrightarrow{\simeq} \sHom(\D_{\mathcal{M}},\D_{\mathcal{M}'})
\]
is an equivalence of categories. On the other hand, we check at once that $u : \mathcal{M} \to \mathcal{M}'$ is left exact (or right exact) if and only if $\D_u : \D_{\mathcal{M}} \to \D_{\mathcal{M}'}$ is so.

Thus, to see whether a morphism of derivators of the form
\[
\begin{tikzcd}[column sep=large, row sep=normal]
f^{\D}_* : &[-2.5em] (\D_{\mathcal{M}})_I \arrow[r] \arrow[d, equal] & (\D_{\mathcal{M}})_{I'} \arrow[d, equal] \\
& \D_{\mathcal{M}(I)} & \D_{\mathcal{M}(I')}
\end{tikzcd}
\]
is \textit{right} exact, it comes to the same thing to say